% !TEX root = ../notes.tex

\documentclass[../notes.tex]{subfiles}

\begin{document}

\section{March 4}
Here we go.

\subsection{The Satake Isomorphism}
Let's continue talking about our parabolic induction. We continue with $G=\op{GL}_n(F)$ with compact open subgroup $K=\op{GL}_n(\OO_F)$.
\begin{remark}
	It turns out that $\op{Ind}_B^G\chi$ always has finite length, and it admits a unique spherical subquotient. In general, permuting the entries of $\chi$ grants the same spherical representation.
\end{remark}
\begin{example}
	For $G=\op{GL}_2(F)$, we saw that $\op{Ind}_B^G\delta^{-1/2}$ admits the subrepresentation $\CC$, but $\op{Ind}_B^G\delta^{1/2}$ admits the quotient $\CC$. This is unsurprising because $\delta^{1/2}$ and $\delta^{-1/2}$ differ merely by a transposition.
\end{example}
Thus, we have constructed a map from $(\CC^\times)^n/S_n$ to the collection of spherical representations.

It remains to show that we have in fact found all the spherical representations.
\begin{remark}
	For any compact open subgroup $K$ of a totally disconnected group $G$, there is a functor $(-)^K\colon\mathrm{Rep}(G)\to\mathrm{Mod}(\mc H_K)$. (In fact, this functor is exact because it passes to finite-dimensional spaces, and fixed points by compact subgroups is exact.) Namely, recall $\mc H_K$ is presented as $e_K\mc H_Ge_K$, which then acts naturally on $V^K=e_KV$.
\end{remark}
We want to make the functor in the previous remark into an equivalence. Before we achieve this, we will find some adjoints.
\begin{definition}
	There is an induction $\op{ind}_!\colon\mathrm{Mod}(\mc H_K)\to\op{Rep}(G)$ given by $M\mapsto\mc He_K\otimes_{\mc H_K}M$. There is similarly an induction $\op{ind}_*$ given by $M\mapsto\op{Hom}_{\mc H_K}(\mc He_K,M)$.
\end{definition}
\begin{remark}
	Both of these functors output to smooth representations.
\end{remark}
\begin{remark}
	One can check that composing either of these functors with $(-)^K$ recovers the original module, so we have ``one-sided inverses.''
\end{remark}
If we would like for either of these functors be equivalences, we need to restrict the category $\op{Rep}_{\mathrm{sm}}(G)$. For example, $V$ should be generated by $V^K$ if it is to come from $V^K$. In fact, this implies that every quotient of $V$ should have $K$-fixed vectors, so we will want to require the ``dual condition'' that every nonzero subrepresentation of $V$ should be $K$-fixed vectors.
\begin{notation}
	Let $G$ be a totally disconnected group, and choose a compact open subgroup $K\subseteq G$. Then $\op{Rep}_{\mathrm{sm}}(G,K)$ is the full subcategory of smooth representations of $G$ such that every quotient and subrepresentation of $V$ admits $K$-fixed vectors.
\end{notation}
\begin{theorem}
	Let $G$ be a totally disconnected group, and choose a compact open subgroup $K\subseteq G$. The functor $(-)^K\colon\op{Rep}_{\mathrm{sm}}(G,K)\to\mathrm{Mod}(\mc H_K)$ is an equivalence.
\end{theorem}
\begin{proof}
	Omitted.
\end{proof}
\begin{corollary}
	Let $G$ be a totally disconnected group, and choose a compact open subgroup $K\subseteq G$. The irreducible smooth representations $V$ of $G$ with $V^K\ne0$ are in bijection with finite-dimensional simple $\mc H_K$-modules.
\end{corollary}
Thus, returning to $\op{GL}_n(F)$, we are interested in understanding the ``spherical'' Hecke algebra
\[\mc H_K=C_c(\op{GL}_n(\OO)\backslash\op{GL}_n(F)/\op{GL}_n(\OO))\,dg.\]
This ring turns out to be commutative (this is a general fact for split reductive groups $G$), and it is not too hard to describe this ring explicitly.
\begin{theorem}[Satake]
	Set $G\coloneqq\op{GL}_n(F)$ and $K\coloneqq\op{GL}_n(\OO)$. Then
	\[\mc H_K\cong\CC\left[t_1^{\pm1},\ldots,t_n^{\pm1}\right]^{S_n}.\]
\end{theorem}
Thus, we see that the simple $\mc H_K$-modules are in bijection with the maximal ideals of $\mc H_K$, which are then in bijection with $\left(\CC^\times\right)/S_n$ by some argument with the Nullstellensatz. It remains to track through the Satake isomorphism to verify that the produced $\mc H_K$-modules do in fact come from our principal series.
\begin{example}
	Let's work with $G=\op{GL}_2(F)$. One may hope that our modules should come from the quotient $\op{GL}_2(F)/\op{GL}_2(\OO)$. This quotient is in bijection with $\OO$-lattices in $F^2$ by sending $g\op{GL}_2(\OO)$ to the lattice $g\OO^2$. If we descend to $\op{PGL}_2(F)/\op{PGL}_2(\OO)$ (which parameterizes lattices up to scalar), there is a bijection to the Bruhat--Tits tree, which is the infinite $(q+1)$-regular tree. Indeed, start with a vertex for the standard lattice $\OO^2$, and then the immediate neighbors are given by those lattices $\OO^2+\varpi^{-1}\ell$, where $\ell\in\PP^1(\FF_q)$ is some line; this ``addition of a line'' edge rule can then be iterated.
\end{example}
\begin{remark}
	Let's explain how this Bruhat--Tits tree can help us: the Hecke algebra can be viewed as $\op{GL}_n(F)$-invariant functions on
	\[\op{GL}_n(F)/\op{GL}_n(\OO)\times\op{GL}_n(F)/\op{GL}_n(\OO).\]
	Thus, for $n=2$, we are basically looking at lattices up to an action by $\op{GL}_2(F)$; it turns out that the only invariant for such pairs is the distance between the two points in the tree!
\end{remark}
\begin{remark}
	For $\op{GL}_n$ in general, this is a Bruhat--Tits building.
\end{remark}

\end{document}